\section{Conceptual Model}\label{sec:conceptual-model}

Live-coding languages often model music as a linear command stream (\texttt{play}, \texttt{sleep}), which suits
performance but makes reusable structure fragile: changing one motif can require rewiring timing throughout the
script.
Motivo instead treats a piece as a \textit{recursive hierarchy}: notes grouped into patterns, patterns sequenced in
tracks, optional voices for parallel lines within a track, so musical content (the ``what'') stays separate from
scheduling (the ``when'').

Patterns behave like pure procedures: they encapsulate logic, accept typed parameters, and produce deterministic
event material when invoked, without a real-time runtime.
That supports modularity (define once, call many times), abstraction (hide polyphonic detail behind a name), and
compile-time resolution of every temporal relationship.

\begin{figure}[htbp]
  \centering
  \motivoscalefig{%
    \begin{tikzpicture}[
        level/.style={motivo/box, minimum width=10.2cm, minimum height=0.9cm},
        arrow/.style={motivo/arrow, line width=1.3pt}
      ]
      \node[level, fill=motivoBlue!15] (program) {Motivo program: tempo, time signature, global definitions};
      \node[level, fill=motivoCyan!18, below=0.35cm of program] (track) {Tracks: instrumental identity and MIDI
      channel allocation};
      \node[level, fill=motivoGreen!18, below=0.35cm of track] (voice) {Voices: parallel timeline fragments
      inside a track};
      \node[level, fill=motivoAmber!20, below=0.35cm of voice] (pattern) {Patterns: reusable parameterized
      musical procedures};
      \node[level, fill=motivoRose!14, below=0.35cm of pattern] (values) {Values: notes, rests, chords,
      sequences, drum notes, scalars, booleans};

      \draw[arrow] (program) -- (track);
      \draw[arrow] (track) -- (voice);
      \draw[arrow] (voice) -- (pattern);
      \draw[arrow] (pattern) -- (values);
    \end{tikzpicture}%
  }
  \caption{Motivo's source-level hierarchy.
  The compiler validates this hierarchy before lowering it into timestamped note events.}\label{fig:motivo-hierarchy}
\end{figure}
